% ============================================================================
% MUSTERLÖSUNG: Kräfte zusammensetzen (UE 2)
% Physik Klasse 7 – Gesamtschule MV
% Version: Kariertes Papier mit TikZ, X-Kreuze
% ============================================================================
\documentclass[a4paper,11pt]{article}

\usepackage[left=1.2cm,right=1.2cm,top=1.5cm,bottom=1.5cm]{geometry}
\usepackage[utf8]{inputenc}
\usepackage[T1]{fontenc}
\usepackage[ngerman]{babel}
\usepackage{tikz}
\usepackage{fancyhdr}
\usepackage{tcolorbox}
\usepackage{amssymb,amsmath}
\usepackage{xcolor}

\renewcommand{\familydefault}{\sfdefault}

\pagestyle{fancy}
\fancyhf{}
\renewcommand{\headrulewidth}{0.4pt}
\lhead{\textbf{\textcolor{red}{MUSTERLÖSUNG}}}
\chead{Kräfte zusammensetzen}
\rhead{Physik Klasse 7}
\setlength{\headheight}{14pt}
\cfoot{\thepage}

\definecolor{physik}{HTML}{2c5aa0}
\definecolor{gridcolor}{RGB}{200,200,200}
\definecolor{axiscolor}{RGB}{100,100,100}
\definecolor{f1color}{RGB}{21,101,192}
\definecolor{f2color}{RGB}{46,125,50}
\definecolor{frescolor}{RGB}{198,40,40}

\newcommand{\luecke}[1]{\underline{\hspace{#1}}}
\newcommand{\lsg}[1]{\textcolor{red}{\textbf{#1}}}

% X-Kreuz (45° gedreht) am Angriffspunkt
\newcommand{\xkreuz}[2]{%
    \draw[black, very thick] (#1-0.4,#2-0.4) -- (#1+0.4,#2+0.4);
    \draw[black, very thick] (#1-0.4,#2+0.4) -- (#1+0.4,#2-0.4);
}

\setlength{\parindent}{0pt}

\begin{document}

% ==================== MERKSATZ-BOX ====================

\begin{tcolorbox}[colback=white,colframe=physik,boxrule=1.5pt,arc=0pt,title={\textbf{Merksatz: Kräfte zusammensetzen}},left=3mm,right=3mm,top=2mm,bottom=2mm]

Ein \textbf{Kraftpfeil} zeigt drei Eigenschaften: \lsg{Angriffspunkt}, \lsg{Richtung}, \lsg{Betrag}

\vspace{2mm}

\begin{minipage}[t]{0.48\textwidth}
\textbf{Gleiche Richtung:}\\[1mm]
$F_{\text{res}} = F_1 + F_2$
\end{minipage}
\hfill
\begin{minipage}[t]{0.48\textwidth}
\textbf{Entgegengesetzt:}\\[1mm]
$F_{\text{res}} = |F_1 - F_2|$
\end{minipage}

\vspace{2mm}

\textbf{Gleichgewicht:} $F_{\text{res}} = 0$ (wenn $F_1 = F_2$ entgegengesetzt) \hfill \textbf{Maßstab:} 1 cm $\hat{=}$ 1 N

\end{tcolorbox}

\vspace{4mm}

% ==================== AUFGABE 1 ====================

\textbf{Aufgabe 1: Kraftpfeile zeichnen}
Zeichne die Kraftpfeile im richtigen Maßstab (1 cm = 1 N = 2 Kästchen). Das $\times$ markiert den Angriffspunkt.

\textbf{Achtung:} Die Pfeilspitze zählt zur Gesamtlänge – nicht zusätzlich draufsetzen!

\vspace{3mm}

\begin{tabular}{@{}p{4.2cm}p{5.5cm}p{5.5cm}@{}}
a) $F = 3$ N nach rechts &
b) $F = 4$ N nach oben &
c) $F = 2{,}5$ N nach links \\[2mm]
\begin{tikzpicture}[scale=0.5]
    \draw[gridcolor, very thin] (0,0) grid (10,4);
    \xkreuz{2}{2}
    \draw[frescolor, very thick, ->] (2,2) -- (8,2);
    \node[frescolor, above, font=\tiny] at (5,2) {3 N};
\end{tikzpicture}
&
\begin{tikzpicture}[scale=0.5]
    \draw[gridcolor, very thin] (0,0) grid (10,10);
    \xkreuz{5}{1}
    \draw[frescolor, very thick, ->] (5,1) -- (5,9);
    \node[frescolor, right, font=\tiny] at (5,5) {4 N};
\end{tikzpicture}
&
\begin{tikzpicture}[scale=0.5]
    \draw[gridcolor, very thin] (0,0) grid (10,4);
    \xkreuz{8}{2}
    \draw[frescolor, very thick, ->] (8,2) -- (3,2);
    \node[frescolor, above, font=\tiny] at (5.5,2) {2,5 N};
\end{tikzpicture}
\end{tabular}

\vspace{6mm}

% ==================== AUFGABE 2 ====================

\textbf{Aufgabe 2: Kräfte in gleicher Richtung}
a) Zeichnung:
\begin{center}
\begin{tikzpicture}[scale=0.5]
    \draw[gridcolor, very thin] (0,0) grid (18,6);
    \draw[axiscolor, thick] (0,3) -- (18,3);
    \xkreuz{2}{3}
    \node[below,font=\small] at (2,2.3) {Wagen};
    % Lösungen
    \draw[f1color, very thick, ->] (2,3.5) -- (8,3.5);
    \node[f1color, above, font=\tiny] at (5,3.5) {$F_1$ = 3 N};
    \draw[f2color, very thick, ->] (2,2.5) -- (10,2.5);
    \node[f2color, below, font=\tiny] at (6,2.5) {$F_2$ = 4 N};
    \draw[frescolor, ultra thick, ->] (2,3) -- (16,3);
    \node[frescolor, above, font=\small] at (9,3) {$F_{\text{res}}$ = 7 N};
\end{tikzpicture}
\end{center}

b) Berechnung: $F_{\text{res}} = $ \lsg{3 N + 4 N} $= $ \lsg{7 N}
\vspace{6mm}

% ==================== AUFGABE 3 ====================

\textbf{Aufgabe 3: Tauziehen}
a) Zeichnung:
\begin{center}
\begin{tikzpicture}[scale=0.5]
    \draw[gridcolor, very thin] (0,0) grid (22,6);
    \draw[axiscolor, thick] (0,3) -- (22,3);
    \xkreuz{13}{3}
    \node[font=\small] at (3,5) {$\leftarrow$ Team A};
    \node[font=\small] at (20,5) {Team B $\rightarrow$};
    % Lösungen (6N = 12 Kästchen links, 4N = 8 Kästchen rechts)
    \draw[f1color, very thick, <-] (1,3.5) -- (13,3.5);
    \node[f1color, above, font=\tiny] at (7,3.5) {$F_1$ = 6 N};
    \draw[f2color, very thick, ->] (13,2.5) -- (21,2.5);
    \node[f2color, below, font=\tiny] at (17,2.5) {$F_2$ = 4 N};
    \draw[frescolor, ultra thick, <-] (9,3) -- (13,3);
    \node[frescolor, below, font=\small] at (11,3) {$F_{\text{res}}$ = 2 N};
\end{tikzpicture}
\end{center}

b) $F_{\text{res}} = $ \lsg{6 N $-$ 4 N} $= $ \lsg{2 N nach links}
c) Welches Team gewinnt? \lsg{Team A}
\vspace{6mm}

% ==================== SEITENUMBRUCH ====================
\newpage

% ==================== AUFGABE 4 ====================

\textbf{Aufgabe 4: Kräftegleichgewicht}
a) $F_{\text{res}} = $ \lsg{5 N $-$ 5 N = 0 N}
\vspace{4mm}

b) Was passiert? \lsg{Nichts, Gleichgewicht (keine Bewegung)}
\vspace{8mm}

% ==================== AUFGABE 5 ====================

\textbf{Aufgabe 5: Kraftpfeil analysieren}
Nenne die drei Eigenschaften eines Kraftpfeils:

\vspace{3mm}

\begin{tabular}{@{}r@{\hspace{3mm}}l@{}}
1. & \lsg{Angriffspunkt} \\[4mm]
2. & \lsg{Richtung} \\[4mm]
3. & \lsg{Betrag / Länge} \\
\end{tabular}

\vspace{8mm}

\hrule
\vspace{3mm}
\textbf{Gesamt: 18 BE}

\end{document}
